% PAGES: 267-267
% First chunk of Chapter 9. Nothing open from Chapter 8: sec08_c2_4.tex (the
% last file of ch08) closed its own note confirming PDF page 267 / folio 251
% opens Chapter 9 fresh, and I independently rendered and read this page,
% confirming the same. The orchestrator writes ch09.tex with the \chapter{}
% and \label{} commands itself (matching the ch01-ch08 pattern), so this file
% does not repeat them.

\noindent Efficient portfolios. Markowitz Portfolio Theory.
\medskip

\noindent Blueprints for finding efficient portfolios.
\medskip

\noindent Minimum variance portfolios. Minimum variance portfolios and the tangency
portfolio.
\medskip

\noindent Maximum return portfolios. Maximum return portfolios and the tangency
portfolio.
\medskip

\noindent Minimum variance portfolio with no cash position.
\medskip

\noindent Value at Risk (VaR). Portfolio VaR.
\medskip

\noindent Subadditivity of VaR and counterexample.
\medskip

\section{Efficient portfolios. Markowitz portfolio theory.}

Consider a portfolio invested in $n$ assets and with a cash position. Let $w_i$ be the
weight of asset $i$ in the portfolio, for $i = 1:n$, i.e., $w_i$ represents the
proportion of the portfolio invested in asset $i$. Then, the weight of the cash
position of the portfolio is
\begin{equation}
w_{cash} \;=\; 1 - \sum_{i=1}^n w_i.
\end{equation}

% Verified at 1200 dpi (folio 251-252): $w$, $\mu$, $\bar\mu$ and their entries are
% set in plain italic throughout this chapter, NOT bold -- unlike the \bfX-style
% vectors/matrices used in earlier chapters. The only symbol actually bolded in
% this chapter is $\mathbf{1}$, the vector of ones (matches the preamble note).
% Reproduced as printed; see also sec09_c2_1.tex's independent confirmation of
% the same convention on later pages of this chapter.
\begin{examplex}
Consider a \$80 million portfolio which is invested in three assets, as follows: a long
\$50 million position in asset 1, a short \$20 million position in asset 2, a long \$25
million position in asset 3, and with a long \$25 million cash position.\footnote{The
asset position in the portfolio is \$50mil + (-\$20mil) + \$25mil = \$55mil, and
therefore the cash position is \$80mil - \$55mil = \$25mil.} The weights $w_1$, $w_2$,
$w_3$, and $w_{cash}$ of the three assets and of the cash position are:
\[
w_1 \;=\; \frac{50}{80} \;=\; 0.625 \;=\; 62.50\%; \qquad w_2 \;=\; -\frac{20}{80} \;=\;
-0.25 \;=\; -25\%;
\]
\[
w_3 \;=\; \frac{25}{80} \;=\; 0.3125 \;=\; 31.25\%; \qquad w_{cash} \;=\; \frac{25}{80}
\;=\; 0.3125 \;=\; 31.25\%.
\]
\end{examplex}
